% !TeX root = main.tex
\chapter{Introducción}
\label{ch:introduccion}

El Sistema Integrado de Predicción y Monitoreo de Sequías Hidrometeorológicas (SIPREH) es una plataforma web diseñada para el análisis continuo, el monitoreo histórico y la predicción probabilística de condiciones de sequía sobre la ciudad de Bogotá D.C.\ y las siete principales cuencas hidrográficas que abastecen su sistema de acueducto: La Regadera, Chisacá, Sisga, Chuza, Neusa, Tominé y San Rafael.

El sistema integra datos climáticos derivados de satélite y productos de reanálisis con salidas de modelos hidrometeorológicos numéricos, presentando los resultados a través de una interfaz geoespacial interactiva orientada a investigadores, gestores de recursos hídricos y tomadores de decisiones públicas.

SIPREH se enmarca en el proyecto \textit{Sistema de Análisis y Alerta de Sequías Meteorológicas e Hidrológicas para Bogotá y para las cuencas hidrográficas y embalses que abastecen de agua a la ciudad}, financiado por la Agencia Distrital para la Educación Superior, la Ciencia y la Tecnología (ATENEA) y desarrollado por el Grupo de Investigación en Ingeniería de los Recursos Hídricos (GIREH) de la Facultad de Ingeniería de la Universidad Nacional de Colombia.

\section{Propósito del sistema}
\label{sec:proposito}

SIPREH fue creado como una herramienta de gestión para apoyar a las entidades responsables de la toma de decisiones en torno a la gestión del recurso hídrico desde una perspectiva de planeación. El sistema permite el estudio de las sequías meteorológicas en la ciudad de Bogotá y sus cuencas abastecedoras mediante un enfoque tanto retrospectivo como predictivo, proporcionando una base técnica e informada para la toma oportuna de decisiones.

Para ello, SIPREH integra información hidrometeorológica, productos derivados del análisis satelital, técnicas de Machine Learning y un tablero interactivo que facilita el acceso, la visualización y la interpretación de la información de manera clara, intuitiva y eficiente.

Sus propósitos centrales son:

\begin{itemize}
    \item Monitorear el estado de la sequía meteorológica en Bogotá y sus cuencas abastecedoras mediante índices de sequía ampliamente utilizados en la literatura científica y en el ámbito institucional.
    \item Analizar series históricas de datos hidrometeorológicos con el fin de consolidar bases de datos robustas que sirvan como insumo para la implementación de técnicas de Machine Learning.
    \item Proveer a las entidades responsables de la gestión del recurso hídrico en Bogotá y sus cuencas abastecedoras predicciones probabilísticas de las condiciones de sequía meteorológica, considerando diferentes horizontes de predicción y niveles de agregación temporal que apoyen la toma de decisiones.
    \item Facilitar la consulta y el análisis temporal de la información mediante la generación de visualizaciones unidimensionales (1D) y bidimensionales (2D), con capacidad de exportación para la elaboración de productos gráficos que favorezcan una interpretación clara e intuitiva de los resultados.
    \item Permitir a la entidad o entidades responsables de la administración del sistema garantizar su continuidad y sostenibilidad mediante la carga, configuración y sincronización de la información que alimenta la visualización y el análisis de datos históricos y predictivos.
\end{itemize}

\section{Objetivos principales del sistema}
\label{sec:objetivos}

El desarrollo de SIPREH persigue los siguientes objetivos técnicos y operativos:

\begin{enumerate}
    \item \textbf{Caracterización retrospectiva y predictiva de la sequía:} Permitir el análisis retrospectivo de la sequía hidrológica y el análisis retrospectivo y predictivo de la sequía meteorológica mediante la consulta integrada de información histórica, horizontes y niveles de agregación para la predicción.

    \item \textbf{Visualización interactiva de la información:} Facilitar la exploración espacial y temporal de la información mediante mapas interactivos, series de tiempo y representaciones gráficas unidimensionales (1D) y bidimensionales (2D), favoreciendo la interpretación de las condiciones de sequía.

    \item \textbf{Consulta y exportación de resultados:} Permitir la consulta y descarga de información histórica y predictiva, así como de gráficos y productos derivados del análisis, para apoyar la toma de decisiones.
    
    \item \textbf{Apoyo a la toma de decisiones:} Proporcionar a las entidades responsables de la gestión del recurso hídrico una herramienta de consulta que facilite el seguimiento de las condiciones de sequía mediante el acceso oportuno a información hidrometeorológica histórica y predictiva.
    
    \item \textbf{Gestión de la plataforma:} Proporcionar herramientas de administración que permitan la configuración del sistema, la sincronización de los conjuntos de datos y la gestión de usuarios.

    \item \textbf{Continuidad y sostenibilidad:} Garantizar la continuidad operativa y la sostenibilidad del sistema mediante mecanismos que faciliten la actualización periódica de la información y la incorporación de nuevas versiones de datos históricos y predictivos.

\end{enumerate}

\section{Público objetivo}
\label{sec:publico}

SIPREH está diseñado para los siguientes perfiles de usuario diferenciados:

\begin{definition}
\textbf{Desarrollador.} Personal encargado de garantizar la continuidad del componente predictivo del sistema mediante el entrenamiento y actualización de modelos de Machine Learning, generando nuevos archivos .parquet con los resultados de las predicciones, los cuales constituyen el insumo para la visualización y el despliegue de la información en SIPREH.
\end{definition}

\begin{definition}
\textbf{Administrador del sistema.} Personal capacitado para la gestión de los archivos .parquet correspondientes a los datos históricos y predictivos utilizados por SIPREH, responsable de la actualización y sincronización de los conjuntos de datos, así como de la administración de usuarios y permisos de acceso al sistema.
\end{definition}

\begin{definition}
\textbf{Tomador de decisiones del recurso hídrico.} Personal perteneciente a entidades como la EAAB, IDIGER e IDEAM, con competencias en la gestión del recurso hídrico, que emplea SIPREH como herramienta de consulta para apoyar la toma de decisiones a partir del análisis retrospectivo de la sequía hidrológica y del análisis retrospectivo y predictivo de la sequía meteorológica.
\end{definition}

\begin{remark}
Las secciones de visualización de datos históricos y de predicción son accesibles para todos los usuarios autenticados. El módulo de administración es exclusivo para usuarios con rol de administrador.
\end{remark}
